\section{Proposed System}
\label{sec:method}

\subsection{Overview}

The system is built as a graph of nodes, each with a single task, in
LangGraph~\cite{langgraph}. The flow of a question through the graph is shown in
Fig.~\ref{fig:pipeline}. The question is first analysed; the planner then selects one
retrieval tool per round, and the results are added to an evidence buffer. The
sufficiency gate (G1) either sends control back to the planner for another round or
releases the buffer to the writer. The writer drafts claims with citations, the
faithfulness gate (G2) checks every claim and returns failed drafts for rewriting, and
the accepted claims are grounded with proof boxes before the answer is returned. Two
loops are therefore present, and both are bounded: at most three retrieval rounds and
at most two rewrites are allowed.

\begin{figure}[!htb]
\centering
\begin{tikzpicture}[
  font=\small,
  node distance=3.2mm,
  box/.style={draw=black!70, rounded corners=2pt, line width=0.6pt, align=center,
              text width=4.6cm, minimum height=6.2mm, inner sep=2pt, fill=#1},
  box/.default=white,
  tool/.style={draw=black!70, rounded corners=2pt, line width=0.6pt, align=center,
               text width=2.0cm, minimum height=6.2mm, inner sep=2pt, fill=orange!12},
  gate/.style={box=red!10, line width=0.9pt},
  arr/.style={-{Stealth[length=2mm]}, line width=0.6pt},
  back/.style={-{Stealth[length=2mm]}, line width=0.6pt, dashed, red!70!black},
  lab/.style={font=\footnotesize, align=left}
]
\node[box=gray!10] (q) {Question};
\node[box=blue!8, below=of q] (ana) {Question analysis: sub-goals and\\ the kind of evidence each needs};
\node[box=violet!10, below=of ana] (plan) {Planner: choose the next tool};
\node[tool, below=4mm of plan] (tab) {Table\\ retriever};
\node[tool, left=1.2mm of tab] (txt) {Text\\ retriever};
\node[tool, right=1.2mm of tab] (fig) {Figure\\ retriever (+ zoom)};
\node[box=teal!10, below=4mm of tab] (buf) {Evidence buffer (linked items)};
\node[gate, below=of buf] (g1) {G1: is the evidence sufficient?};
\node[box=orange!18, below=5mm of g1] (wr) {Writer: claims, each citing a\\ quote, a table cell or a figure};
\node[gate, below=of wr] (g2) {G2: is every claim supported?};
\node[box=green!12, below=of g2] (out) {Proof boxes, claim links and\\ trace $\rightarrow$ answer};

\draw[arr] (q) -- (ana);
\draw[arr] (ana) -- (plan);
\draw[arr] (plan) -- (tab);
\draw[arr] (plan.south) -- ++(0,-1.5mm) -| (txt.north);
\draw[arr] (plan.south) -- ++(0,-1.5mm) -| (fig.north);
\draw[arr] (tab) -- (buf);
\draw[arr] (txt.south) -- ++(0,-1.5mm) -| (buf.north);
\draw[arr] (fig.south) -- ++(0,-1.5mm) -| (buf.north);
\draw[arr] (buf) -- (g1);
\draw[arr] (g1) -- node[lab, right=1mm] {yes, or best effort (flagged)} (wr);
\draw[arr] (wr) -- (g2);
\draw[arr] (g2) -- node[lab, right=1mm] {yes} (out);
% the two loops run outside the row of tools, with their labels on the horizontal segment
\coordinate (xl) at (-3.75,0);
\coordinate (xr) at (3.75,0);
\draw[back] (g1.west) -- node[lab, above, midway] {no ($\leq$3)} (g1.west -| xl) |- (plan.west);
\draw[back] (g2.east) -- node[lab, above, midway] {no ($\leq$2)} (g2.east -| xr) |- (wr.east);
\end{tikzpicture}
\caption{Flow of a question through the system. Solid arrows show the forward path.
Dashed arrows show the two bounded loops: the sufficiency gate (G1) returns control to
the planner for another round (at most three), and the faithfulness gate (G2) returns
the draft to the writer (at most two rewrites).}
\label{fig:pipeline}
\end{figure}

\subsection{Question Analysis and Planning}

The question is split into sub-goals, and the names it mentions (datasets, methods,
metrics) are extracted for later filtering. The kind of evidence each sub-goal needs,
a table or a figure, is not left to the language model. It is decided by a vote of the
nine most similar labelled SPIQA questions taken from papers that are used in neither
the development nor the test split. On the development questions, this vote chose the
reference's kind in 17 of 20 cases, whereas the model alone chose it in 8. The planner
then selects one tool per round and writes the query for it; a tool that has already
been used for a sub-goal may be called again only with a new query.

\subsection{Retrieval Tools}

Three tools can be selected by the planner. The text retriever ranks the paper's
paragraphs by reciprocal rank fusion~\cite{cormack2009rrf} of two lists: SPECTER2
embeddings~\cite{singh2023scirepeval} searched with FAISS~\cite{johnson2019faiss}, and
keyword overlap weighted by inverse document frequency. The keyword list was added
because, within one paper, every paragraph shares the topic that the embedding
captures, whereas SPIQA questions reuse the paper's own terms. The table retriever
ranks the paper's tables by their caption words and by the visual score of their page,
and returns each table as cells recovered with Table Transformer~\cite{smock2022pubtables},
each cell keeping its page coordinates. The figure retriever scores page images with
ColQwen2, a model of the ColPali family~\cite{faysse2025colpali}, and returns the best
figure. A fourth tool, a zoom that re-renders the figure at 300~dpi, is not chosen by
the planner; it is run automatically when a figure covers less than 30\% of its page.
Every retrieved item enters the evidence buffer with its page, tool and round, and
items are linked when they share a page or when a sentence names a table or figure.

\subsection{Sufficiency Gate}

Before anything is written, the buffer is judged by G1. The gate is asked to name the
facts in the buffer that answer the question, with their evidence identifiers, and to
give its verdict only afterwards. The order matters because the reply is generated as
JSON in schema order: when the verdict came first, it was justified only after it had
been made. If G1 is not convinced after three rounds, two cases are separated. When no
answering fact was found, the system abstains and reports what is missing. When such
facts were found but judged incomplete, the answer is written from them as a
\emph{best effort} and is marked as possibly incomplete. In this way, G1 decides when
the search stops, while the decision to publish a claim is left to G2.

\subsection{Writing and Faithfulness Gate}

The writer receives the buffer as text, cells and figure images, and returns a list of
claims. Every claim must cite its evidence as a verbatim quote, a table cell given by
row, column and value, or an observation about a figure. Each citation is verified by
G2 in two layers. First, a quote must occur in the cited paragraph and a cell value
must equal the parsed cell; both sides are normalised beforehand for the artefacts of
PDF text extraction, such as ligatures, line-end hyphens and spaced decimals. Second,
the model is asked whether the cited evidence supports the claim as written. A failed
draft is returned with the reason, and claims that still fail after two rewrites are
withheld and listed as unconfirmed.

\subsection{Grounding the Claims}

Every published claim is given a box on the page, and exact sources are preferred: the
cited table cell, the bar of a vector chart read directly from the PDF, the quoted
sentence located on the page, the caption, or the labels of a diagram that the claim
names. Only for photographs, which carry no printed words, is the box taken from the
highest-scoring patches of the ColQwen2 similarity map computed for the claim's text.
An execution trace, listing the tools called and every gate decision, is returned with
the answer.
